\scanpage{34}
écriture intrinsèque, de façon que l'opération de $\Gamma\unsure{\Rightarrow|=|\simeq}\deleted{\operatorname{Aut}}(\Pi)$ sur $E(\Pi)$ sera obtenue automatiquement. (Il restera alors simplement à décrire $S\longrightarrow E(\Pi)$ compatible avec l'op. de $\Gamma$\ldots).

$\overline{\underline{\mathrm{VI}}}$ \uline{Opérations} $\psi_r$ \uline{d'étoilage d'un polygône.}

Soit $(S,\omega)$ un polygône combinatoire, et soit $r\in\mathbb Z$, considérons
\[
\psi_r(\omega)=\{u^r\mid u\in\omega\}\tag*{(23)}
\]
\deleted{\unsure{S'il y a|A' que}}

\uline{Prop 6} a) Pour que \deleted{$\unsure{\psi_r(\omega)}$} $(S,\psi_r(\omega))$ soit un polygône combinatoire (i.e. pour que $u^r$ opère transitivement sur $S$) il f. et s. que \deleted{$S$} $\operatorname{Card}(S)=n$ soit fini, et que $(r,n)=1$. \deleted{Si $\operatorname{card}(S)=n$,}

\inserted{b) On a, pour $r,r'\in\mathbb Z$}
\[
\inserted{\psi_r(\psi_{r'}(\omega))=\psi_{rr'}(\omega)}
\]

c) Si \deleted{$\operatorname{card}S=$} $n$ est fini, alors \deleted{\unsure{l'on a}} \deleted{\unsure{pour les} $\psi_r(\omega)$} et \unsure{soit|pour} $r,r'\in\mathbb Z$, \deleted{\unsure{on a 1)}} alors on a
\[
\psi_r(\omega)=\psi_{r'}(\omega)\quad\text{ssi}\quad r'\equiv\pm r\ (n)
\]
\inserted{donc $\Psi_n$} \deleted{\unsure{donne}} \deleted{$r\mapsto\psi_r(\omega)$ \unsure{induit une application injective} $\unsure{\to}$} \deleted{$\mathbb Z_n^*/\pm\longrightarrow\operatorname{Pol}(S)$} \inserted{opère librement sur $\operatorname{Pol}(S)$}

i.e. \deleted{\unsure{que}} $(\mathbb Z/n\mathbb Z)^*$ \uline{opère} sur l'ens $\operatorname{Pol}(S)$ des structures polygonales sur $S$. En fait, c'est
% margin note written sideways
\margin{(24) $(\mathbb Z/n\mathbb Z)^*/\{\pm1\}\mathrel{\underset{\mathrm{déf}}{=}:}\Psi_n$ qui opère sur $\operatorname{Pol}(S)$, car $-1$ opère de façon triviale}
% archive sheet number: 33
